% !TeX root = ../../../main.tex

\begin{table}[tbp]
    \centering
    \caption{Model Definitions and Assessment Approaches for Regression and Classification Models Predicting Clinical Variables from VisR-Derived Metrics}
    \begin{tabular}{llrll}
        \toprule
        \multirow{2}{*}{Parameter} & 
        \multicolumn{2}{c}{Regressor} & 
        \multicolumn{2}{c}{Classifier}\\
         & Metric & Baseline & Model Type & Categories\\
        \midrule
        Serum Creatinine (mg/dl) &
        RMSE & 0.447 &
        Binary & $<$ 1.6, $\geq$ 1.6\\
        UPCR (g/g) &
        RMSE & 0.589 &
        Binary & $<$ 0.862, $\geq$ 0.150\\
        DSA Presence &
        - & - &
        Binary & Absent, Present\\
        Indication Biopsy &
        - & - &
        Binary & No, Yes\\
        \emph{i} Score &
        - & - &
        Binary & 0, 1--3\\
        IFTA Grade &
        Cross-Entropy & 1.39 &
        Multinomial & 0, 1, 2, 3\\
        \emph{i-IFTA} Score &
        Cross-Entropy & 1.10 &
        Multinomial & 0, 1, 2\\
        Rejection &
        - & - &
        Binary & No, Yes\\
        \bottomrule
    \end{tabular}
    \par\vspace{0.5em}
    \noindent{}RMSE: root mean square error. Baseline metrics for RMSE are calculated as the root mean square of the difference between individual values and their medians. Baseline metrics for cross-entropy losses are based on naive models that randomly predict each category with equal probability. For example, if a clinical parameter has three categories, the baseline is $-ln \left( 1 / 3 \right) \approx 1.10$.
    \label{tbl:clinicalVariables}
\end{table}
